% Einheit 3 von 5 – Extrem- und Wendepunkte mit Parameter (bank/funktionsscharen-und-ortskurven/e3.jsonl, zusammenbau v0.1)
%% TODO Zeitmarke Einheit 3: Marken-Zeile nicht lesbar
%% TODO Zweigzeile Teil 1: was hier gelernt wird (Satz in Schülersprache aus dem Einheitstitel) – Einheit 3
\einheitenkopf[e3]{Einheit 3 von 5 · Extrem- und Wendepunkte mit Parameter}
\zweigzeile{Abitur LK}
\setcounter{aufgabe}{13}

%% TODO Titel: Ich-kann-Satz fehlt in der Bank (Platzhalter: Kettenname) – Nr. 14
%% TODO Anweisung: ein Satz über der ersten Teilaufgabe fehlt in der Bank – Nr. 14
\begin{aufgabe}{Extrem- und Wendepunkte mit Parameter}
%% TODO \rechenplatz{n}: Schrittzahl des Grundfalls fehlt in der Bank (nur bei fester Schreibform, 2.3 b)
\begin{teile}
\teil Gegeben ist $f_a(x) = a x^3 - 2x^2$ mit $a \in \mathbb{R}$. Bilde $f_a'(x)$ und berechne $f_a'(1)$. $f_a'(x) =$ \_\_, $f_a'(1) =$ \_\_
\teil Gegeben ist $f_k(x) = -k \cdot (x^4 - 8x^3)$ mit $k > 0$. Bestimme die x-Koordinate des Hochpunkts. \hfill (Abitur 2020 LK) \\ $x =$ \_\_
\teil Der Punkt $P(1 \,|\, -0{,}5)$ liegt auf dem Graphen von $f_a(x) = \frac{1}{2}x^3 + a x^2 + a x + 2$. Bestimme $a$ und weise nach, dass $P$ Wendepunkt ist. \hfill (Abitur 2026 LK) \\ $a =$ \_\_
\teil Gegeben ist $f_a(x) = \mathrm{ln}(a x^2 + 4)$ mit $a > 0$. Weise nach, dass alle Graphen einen gemeinsamen Extrempunkt haben, und begründe ohne zweite Ableitung, dass er ein Tiefpunkt ist. \hfill (Abitur 2017 LK)
\teil Weise nach, dass die Graphen von $f_a(x) = (x^2 + a) \cdot e^{0{,}5x}$ für $a > 4$ keine Extrempunkte haben. \hfill (Abitur 2018 LK)
\teil Gegeben ist $K_b(x) = x^3 - b x^2 + 12x + 5$ mit $b > 0$. Weise nach, dass $K_b$ für $b < 6$ streng monoton wachsend ist. \hfill (Abitur 2018 LK)
\teil Gegeben ist $f_k(x) = \frac{1}{k} \cdot x^2 \cdot (x - 4k)^2$ mit $k > 0$. Die Tiefpunkte liegen bei $x = 0$ und $x = 4k$, der Hochpunkt bei $x = 2k$. Berechne den Abstand des Hochpunkts zu jedem Tiefpunkt in Abhängigkeit von $k$. \hfill (Abitur 2025 LK)
\teil Gegeben ist $f_a(x) = \frac{1}{a}x^3 + 2x^2 + 3x - a$ mit $a \neq 0$. Für genau einen Wert von $a$ hat der Graph genau einen Punkt mit waagerechter Tangente. Bestimme ihn und erläutere, wie man nachweisen könnte, dass dort ein Sattelpunkt liegt. \hfill (Abitur 2018 LK)
\end{teile}
\end{aufgabe}

%% TODO Titel: Ich-kann-Satz fehlt in der Bank (Platzhalter: Kettenname) – Nr. 15
%% TODO Anweisung: ein Satz über der ersten Teilaufgabe fehlt in der Bank – Nr. 15
\begin{aufgabe}{Extrem- und Wendepunkte mit Parameter – Fehler finden, Begründen, Anwendung}
\begin{teile}
\teil Ben bestimmt die Extrempunkte von $f_a(x) = x^3 - 3a x$ mit $a > 0$: „$f_a'(x) = 3x^2 - 3a = 0$, also Hochpunkt bei $x = \sqrt{a}$ und Tiefpunkt bei $x = -\sqrt{a}$.“ Finde den Fehler und korrigiere.
\teil Begründe ohne vollständige Rechnung, warum der Graph von $f_a(x) = x^3 + a x$ für $a > 0$ keinen Extrempunkt hat.
\teil Die Temperatur einer Flüssigkeit wird durch $f_k(t) = 20 + 30t \cdot e^{-kt}$ mit $k > 0$ beschrieben ($t$ in min, $f_k(t)$ in °C). Bestimme in Abhängigkeit von $k$, wann die Temperatur am höchsten ist und wann sie am stärksten abnimmt. \hfill (Abitur 2017 GK)
\end{teile}
\end{aufgabe}
